% Options for packages loaded elsewhere
\PassOptionsToPackage{unicode}{hyperref}
\PassOptionsToPackage{hyphens}{url}
\documentclass[
]{article}
\usepackage{xcolor}
\usepackage[margin=1in]{geometry}
\usepackage{amsmath,amssymb}
\setcounter{secnumdepth}{-\maxdimen} % remove section numbering
\usepackage{iftex}
\ifPDFTeX
  \usepackage[T1]{fontenc}
  \usepackage[utf8]{inputenc}
  \usepackage{textcomp} % provide euro and other symbols
\else % if luatex or xetex
  \usepackage{unicode-math} % this also loads fontspec
  \defaultfontfeatures{Scale=MatchLowercase}
  \defaultfontfeatures[\rmfamily]{Ligatures=TeX,Scale=1}
\fi
\usepackage{lmodern}
\ifPDFTeX\else
  % xetex/luatex font selection
\fi
% Use upquote if available, for straight quotes in verbatim environments
\IfFileExists{upquote.sty}{\usepackage{upquote}}{}
\IfFileExists{microtype.sty}{% use microtype if available
  \usepackage[]{microtype}
  \UseMicrotypeSet[protrusion]{basicmath} % disable protrusion for tt fonts
}{}
\makeatletter
\@ifundefined{KOMAClassName}{% if non-KOMA class
  \IfFileExists{parskip.sty}{%
    \usepackage{parskip}
  }{% else
    \setlength{\parindent}{0pt}
    \setlength{\parskip}{6pt plus 2pt minus 1pt}}
}{% if KOMA class
  \KOMAoptions{parskip=half}}
\makeatother
\usepackage{color}
\usepackage{fancyvrb}
\newcommand{\VerbBar}{|}
\newcommand{\VERB}{\Verb[commandchars=\\\{\}]}
\DefineVerbatimEnvironment{Highlighting}{Verbatim}{commandchars=\\\{\}}
% Add ',fontsize=\small' for more characters per line
\usepackage{framed}
\definecolor{shadecolor}{RGB}{248,248,248}
\newenvironment{Shaded}{\begin{snugshade}}{\end{snugshade}}
\newcommand{\AlertTok}[1]{\textcolor[rgb]{0.94,0.16,0.16}{#1}}
\newcommand{\AnnotationTok}[1]{\textcolor[rgb]{0.56,0.35,0.01}{\textbf{\textit{#1}}}}
\newcommand{\AttributeTok}[1]{\textcolor[rgb]{0.13,0.29,0.53}{#1}}
\newcommand{\BaseNTok}[1]{\textcolor[rgb]{0.00,0.00,0.81}{#1}}
\newcommand{\BuiltInTok}[1]{#1}
\newcommand{\CharTok}[1]{\textcolor[rgb]{0.31,0.60,0.02}{#1}}
\newcommand{\CommentTok}[1]{\textcolor[rgb]{0.56,0.35,0.01}{\textit{#1}}}
\newcommand{\CommentVarTok}[1]{\textcolor[rgb]{0.56,0.35,0.01}{\textbf{\textit{#1}}}}
\newcommand{\ConstantTok}[1]{\textcolor[rgb]{0.56,0.35,0.01}{#1}}
\newcommand{\ControlFlowTok}[1]{\textcolor[rgb]{0.13,0.29,0.53}{\textbf{#1}}}
\newcommand{\DataTypeTok}[1]{\textcolor[rgb]{0.13,0.29,0.53}{#1}}
\newcommand{\DecValTok}[1]{\textcolor[rgb]{0.00,0.00,0.81}{#1}}
\newcommand{\DocumentationTok}[1]{\textcolor[rgb]{0.56,0.35,0.01}{\textbf{\textit{#1}}}}
\newcommand{\ErrorTok}[1]{\textcolor[rgb]{0.64,0.00,0.00}{\textbf{#1}}}
\newcommand{\ExtensionTok}[1]{#1}
\newcommand{\FloatTok}[1]{\textcolor[rgb]{0.00,0.00,0.81}{#1}}
\newcommand{\FunctionTok}[1]{\textcolor[rgb]{0.13,0.29,0.53}{\textbf{#1}}}
\newcommand{\ImportTok}[1]{#1}
\newcommand{\InformationTok}[1]{\textcolor[rgb]{0.56,0.35,0.01}{\textbf{\textit{#1}}}}
\newcommand{\KeywordTok}[1]{\textcolor[rgb]{0.13,0.29,0.53}{\textbf{#1}}}
\newcommand{\NormalTok}[1]{#1}
\newcommand{\OperatorTok}[1]{\textcolor[rgb]{0.81,0.36,0.00}{\textbf{#1}}}
\newcommand{\OtherTok}[1]{\textcolor[rgb]{0.56,0.35,0.01}{#1}}
\newcommand{\PreprocessorTok}[1]{\textcolor[rgb]{0.56,0.35,0.01}{\textit{#1}}}
\newcommand{\RegionMarkerTok}[1]{#1}
\newcommand{\SpecialCharTok}[1]{\textcolor[rgb]{0.81,0.36,0.00}{\textbf{#1}}}
\newcommand{\SpecialStringTok}[1]{\textcolor[rgb]{0.31,0.60,0.02}{#1}}
\newcommand{\StringTok}[1]{\textcolor[rgb]{0.31,0.60,0.02}{#1}}
\newcommand{\VariableTok}[1]{\textcolor[rgb]{0.00,0.00,0.00}{#1}}
\newcommand{\VerbatimStringTok}[1]{\textcolor[rgb]{0.31,0.60,0.02}{#1}}
\newcommand{\WarningTok}[1]{\textcolor[rgb]{0.56,0.35,0.01}{\textbf{\textit{#1}}}}
\usepackage{graphicx}
\makeatletter
\newsavebox\pandoc@box
\newcommand*\pandocbounded[1]{% scales image to fit in text height/width
  \sbox\pandoc@box{#1}%
  \Gscale@div\@tempa{\textheight}{\dimexpr\ht\pandoc@box+\dp\pandoc@box\relax}%
  \Gscale@div\@tempb{\linewidth}{\wd\pandoc@box}%
  \ifdim\@tempb\p@<\@tempa\p@\let\@tempa\@tempb\fi% select the smaller of both
  \ifdim\@tempa\p@<\p@\scalebox{\@tempa}{\usebox\pandoc@box}%
  \else\usebox{\pandoc@box}%
  \fi%
}
% Set default figure placement to htbp
\def\fps@figure{htbp}
\makeatother
\setlength{\emergencystretch}{3em} % prevent overfull lines
\providecommand{\tightlist}{%
  \setlength{\itemsep}{0pt}\setlength{\parskip}{0pt}}
\usepackage{bookmark}
\IfFileExists{xurl.sty}{\usepackage{xurl}}{} % add URL line breaks if available
\urlstyle{same}
\hypersetup{
  pdftitle={Leave-One-Out Cross-Validation},
  hidelinks,
  pdfcreator={LaTeX via pandoc}}

\title{Leave-One-Out Cross-Validation}
\author{}
\date{\vspace{-2.5em}2026-04-17}

\begin{document}
\maketitle

\subsection{Introduction}\label{introduction}

Leave-one-out cross-validation (LOO-CV) is the gold standard for
estimating out-of-sample predictive performance, but exact LOO requires
refitting the model once per observation --- prohibitive for any
Bayesian model that takes more than a few seconds to sample.
Pareto-Smoothed Importance Sampling LOO (PSIS-LOO), introduced by
Vehtari, Gelman, and Gabry, approximates exact LOO from a single
full-data posterior using importance weights, while providing
per-observation diagnostics that flag where the approximation breaks
down. It has become the standard model-comparison tool for
\texttt{brms}, \texttt{rstanarm}, and raw Stan workflows.

\subsection{Prerequisites}\label{prerequisites}

A working understanding of cross-validation, importance sampling, and
the role of the log pointwise predictive density (elpd) in measuring
predictive performance.

\subsection{Theory}\label{theory}

The expected log predictive density under leave-one-out is

\[\mathrm{elpd}_{\mathrm{loo}} = \sum_{i=1}^n \log p(y_i \mid y_{-i}),\]

where \(y_{-i}\) is the data with observation \(i\) removed. PSIS
approximates each \(p(y_i \mid y_{-i})\) via importance weights
\(w_i^{(s)} \propto 1 / p(y_i \mid \theta^{(s)})\) applied to the
full-data posterior draws \(\theta^{(s)}\). To stabilise the right tail
of the importance weights, PSIS fits a generalised Pareto distribution
to the largest weights and replaces them with order-statistic estimates.
The shape parameter \(\hat k\) of that fit acts as a per-observation
diagnostic: \(\hat k < 0.5\) is excellent, \(0.5 \le \hat k < 0.7\)
acceptable, and \(\hat k > 0.7\) indicates that PSIS-LOO is unreliable
for that point.

\subsection{Assumptions}\label{assumptions}

Exchangeable observations and a posterior with sufficient effective
sample size to estimate the importance-weight tail; for hierarchical or
temporal data, replace standard LOO with leave-one-cluster-out or
leave-future-out variants whose theoretical properties match the
dependence structure.

\subsection{R Implementation}\label{r-implementation}

\begin{Shaded}
\begin{Highlighting}[]
\FunctionTok{library}\NormalTok{(loo); }\FunctionTok{library}\NormalTok{(brms)}

\FunctionTok{data}\NormalTok{(mtcars)}
\NormalTok{fit1 }\OtherTok{\textless{}{-}} \FunctionTok{brm}\NormalTok{(mpg }\SpecialCharTok{\textasciitilde{}}\NormalTok{ wt, }\AttributeTok{data =}\NormalTok{ mtcars, }\AttributeTok{chains =} \DecValTok{2}\NormalTok{, }\AttributeTok{iter =} \DecValTok{2000}\NormalTok{, }\AttributeTok{refresh =} \DecValTok{0}\NormalTok{)}
\NormalTok{fit2 }\OtherTok{\textless{}{-}} \FunctionTok{brm}\NormalTok{(mpg }\SpecialCharTok{\textasciitilde{}}\NormalTok{ wt }\SpecialCharTok{+}\NormalTok{ hp, }\AttributeTok{data =}\NormalTok{ mtcars, }\AttributeTok{chains =} \DecValTok{2}\NormalTok{, }\AttributeTok{iter =} \DecValTok{2000}\NormalTok{, }\AttributeTok{refresh =} \DecValTok{0}\NormalTok{)}

\NormalTok{loo1 }\OtherTok{\textless{}{-}} \FunctionTok{loo}\NormalTok{(fit1)}
\NormalTok{loo2 }\OtherTok{\textless{}{-}} \FunctionTok{loo}\NormalTok{(fit2)}

\FunctionTok{loo\_compare}\NormalTok{(loo1, loo2)}

\CommentTok{\# k{-}hat diagnostic}
\FunctionTok{plot}\NormalTok{(loo2)}
\end{Highlighting}
\end{Shaded}

\subsection{Output \& Results}\label{output-results}

\texttt{loo()} returns the elpd, the effective parameter count
\(p_{\mathrm{loo}}\), and a vector of per-observation \(\hat k\) values.
\texttt{loo\_compare()} aligns two \texttt{loo} objects and returns the
elpd difference with its standard error, ordered from best to worst.
\texttt{plot(loo)} displays \(\hat k\) values against observation index,
immediately revealing influential points.

\subsection{Interpretation}\label{interpretation}

A reporting sentence: ``PSIS-LOO favoured the two-predictor model with
\(\Delta\mathrm{elpd} = 6\) (SE 2.3); all \(\hat k < 0.5\), so the
importance-sampling approximation is reliable, and no observation needs
an exact refit.'' When some \(\hat k\) values exceed 0.7, rerun with
\texttt{loo(fit,\ reloo\ =\ TRUE)} to refit the model on the affected
leave-out folds; the elpd estimate is then a hybrid of PSIS for safe
folds and exact fits for the troublesome ones.

\subsection{Practical Tips}\label{practical-tips}

\begin{itemize}
\tightlist
\item
  Check \(\hat k\) values \emph{before} interpreting an elpd comparison;
  a single high-leverage observation can dominate the ranking.
\item
  For time-series, use leave-future-out CV via
  \texttt{loo::loo\_subsample()} or \texttt{loo\_moment\_match()};
  standard LOO leaks future information into past folds and biases elpd
  downward.
\item
  For hierarchical models with few clusters, leave-one-cluster-out CV is
  the appropriate generalisation: it answers the prediction question
  that matters in practice (``how would the model perform on a new
  cluster?'').
\item
  Differences in elpd are interpretable in absolute terms:
  \(\Delta\mathrm{elpd} > 4 \times \mathrm{SE}\) is decisive;
  \(|\Delta\mathrm{elpd}| < 2 \times \mathrm{SE}\) is inconclusive.
\item
  LOO is well-defined for non-nested models, unlike likelihood-ratio
  tests, which makes it the natural tool when comparing models that
  differ in family, link, or predictor set.
\end{itemize}

\subsection{R Packages Used}\label{r-packages-used}

\texttt{loo} for \texttt{loo()}, \texttt{loo\_compare()},
\texttt{loo\_subsample()}, and \texttt{loo\_moment\_match()},
\texttt{brms} and \texttt{rstanarm} for fits that expose pointwise
log-likelihoods automatically, and \texttt{bayesplot} for posterior
predictive plots that complement the elpd-based comparison.

\subsection{For Reviewers}\label{for-reviewers}

\textbf{What to look for in a paper using this method.}

\begin{itemize}
\tightlist
\item
  \textbf{Common misapplications.}
\item
  \textbf{Diagnostics that should be reported but often aren't.}
\item
  \textbf{Red flags in tables and figures.}
\item
  \textbf{What to verify.}
\item
  \textbf{What an adequate Methods paragraph must contain.}
\end{itemize}

\subsection{See also --- labs in this
chapter}\label{see-also-labs-in-this-chapter}

\begin{itemize}
\tightlist
\item
  \href{labs/lab-bayesian-thinking-control-vs-decision-error.Rmd}{Bayesian
  thinking; α control vs decision error}
\item
  \href{labs/lab-brms-stan-loo-hierarchical-models.Rmd}{brms / Stan,
  LOO, hierarchical models}
\end{itemize}

\end{document}
